% ======================================================================
% Раздел 1. Двойные интегралы (задачи 1.1 и 1.2).
% Ответственный: D — Филиппов Андрей.
%
% Что сдать (main.tex: порядок 249–283, разобранный пример 288–387;
% требования к отчёту 1854–1861; правила точности 1948–1999;
% пример таблицы и графика ошибки 2001–2083):
%   1.1  Аналитически вычислить I1 по прямоугольнику [0,a]x[0,b];
%        реализовать метод средних прямоугольников и метод трапеций;
%        вычислить I1 при n = 4, 8, 16, 32, 64 по каждой координате.
%   1.2  Аналитически вычислить I2 по треугольнику D; перейти к единичному
%        квадрату (x=a*xi, y=b*(1-xi)*eta), найти якобиан и вычислить
%        преобразованный интеграл численно обоими методами.
%   Для обеих задач: изображения областей, таблицы сравнения с аналитикой,
%   графики абсолютной ошибки от шага, вывод об изменении ошибки при удвоении n.
%   Нулевые ошибки на логарифмическую шкалу не наносить; если метод даёт точное
%   значение с учётом округления — указать это в отчёте.
%
% Параметры варианта 3 (приложение А): a=2, b=3.
% ВНИМАНИЕ: a,b,c приложения А (интегралы) и приложения Б (модель) — разные величины.
% Шаги: h_x=a/n, h_y=b/n; на графике h=max(h_x,h_y); после замены h=1/n.
%
% Код квадратур — в solution.py (блок D), вызовы — task_1_1/task_1_2
% в experiments.py (запуск: python experiments.py 1.1).
% Рисунки: figures/1-1-*.pdf, figures/1-2-*.pdf (сохранять из experiments.py).
% ======================================================================

\subsection{Интеграл по прямоугольной области}

\taskhead{Исходные данные и аналитическое решение}
Для варианта 3 заданы $a=2$, $b=3$. Рассматривается интеграл
\[
I_1=\int_0^2\int_0^3(x^2+xy+y^2)\,dy\,dx.
\]
\todo{D: вычислить $I_1$ аналитически — повторное интегрирование с промежуточными шагами (пример: main.tex, строки 292–304).}

\taskhead{Программный результат и сравнение}
\todo{D: листинг вызова своих функций из solution.py и полученные значения (образец: \texttt{report\_template.tex}, строки 138–149).}
Аналитическое значение $I_1$: \todo{D: подставить}. Шаги сетки $h_x=a/n$, $h_y=b/n$;
для графика используется $h=\max(h_x,h_y)$. Ошибки $\Delta_n=|I_n-I_1|$,
$\delta_n=\Delta_n/|I_1|$ вычисляются до округления результатов.
\begin{table}[htbp]
\centering\small
\caption{Сравнение численного интеграла с аналитическим значением $I_1$}
\label{tab:1-1}
\begin{tabular}{|l|r|r|r|r|r|}\hline
Метод & $n$ & $h$ & $I_n$ & $\Delta_n$ & $\delta_n$\\\hline
Ср. прямоугольники & 4 & — & — & — & —\\
Ср. прямоугольники & 8 & — & — & — & —\\
Ср. прямоугольники & 16 & — & — & — & —\\
Ср. прямоугольники & 32 & — & — & — & —\\
Ср. прямоугольники & 64 & — & — & — & —\\
Трапеции & 4 & — & — & — & —\\
Трапеции & 8 & — & — & — & —\\
Трапеции & 16 & — & — & — & —\\
Трапеции & 32 & — & — & — & —\\
Трапеции & 64 & — & — & — & —\\
\hline\end{tabular}
\end{table}

\taskhead{Графическая проверка}
\taskfigure[0.6]{figures/1-1-region.pdf}{Область интегрирования — прямоугольник $[0,a]\times[0,b]$ (равный масштаб осей)}{fig:1-1-region}
\taskfigure[0.75]{figures/1-1-error.pdf}{Абсолютная ошибка обоих методов от шага сетки (логарифмические оси)}{fig:1-1-error}
\todo{D: пояснить графики: как меняется ошибка при удвоении $n$ и почему.}

\taskhead{Вывод}
\todo{D: вывод о сходимости обоих методов; если метод даёт точное значение с учётом округления — указать это.}

\subsection{Интеграл по треугольной области}

\taskhead{Исходные данные и аналитическое решение}
Для варианта 3 заданы $a=2$, $b=3$. Область
$D=\{x\geq0,\ y\geq0,\ x/a+y/b\leq1\}$ --- треугольник с вершинами $(0,0)$, $(a,0)$, $(0,b)$.
\[
I_2=\iint_D(x+y)\,dx\,dy.
\]
\todo{D: выписать пределы интегрирования для своего варианта и вычислить $I_2$ аналитически (пример: main.tex, строки 348–357).}

\taskhead{Замена переменных и якобиан}
% Переход к единичному квадрату: x=a*xi, y=b*(1-xi)*eta, 0<=xi,eta<=1
% (теория: main.tex, строки 225–247; пример: 359–375).
\todo{D: выписать замену для своего варианта, матрицу частных производных, вычислить якобиан $J$; записать преобразованный интеграл по единичному квадрату от $g(\xi,\eta)=f\bigl(x(\xi,\eta),y(\xi,\eta)\bigr)\,|J|$.}

\taskhead{Программный результат и сравнение}
\todo{D: листинг вызовов функций; оба метода применяются к $g$ на единичном квадрате, $h_\xi=h_\eta=1/n$.}
Аналитическое значение $I_2$: \todo{D: подставить}.
\begin{table}[htbp]
\centering\small
\caption{Сравнение численного интеграла с аналитическим значением $I_2$}
\label{tab:1-2}
\begin{tabular}{|l|r|r|r|r|r|}\hline
Метод & $n$ & $h$ & $I_n$ & $\Delta_n$ & $\delta_n$\\\hline
Ср. прямоугольники & 4 & — & — & — & —\\
Ср. прямоугольники & 8 & — & — & — & —\\
Ср. прямоугольники & 16 & — & — & — & —\\
Ср. прямоугольники & 32 & — & — & — & —\\
Ср. прямоугольники & 64 & — & — & — & —\\
Трапеции & 4 & — & — & — & —\\
Трапеции & 8 & — & — & — & —\\
Трапеции & 16 & — & — & — & —\\
Трапеции & 32 & — & — & — & —\\
Трапеции & 64 & — & — & — & —\\
\hline\end{tabular}
\end{table}

\taskhead{Графическая проверка}
\taskfigure[0.6]{figures/1-2-region.pdf}{Область интегрирования — треугольник $D$ (равный масштаб осей)}{fig:1-2-region}
\taskfigure[0.75]{figures/1-2-error.pdf}{Абсолютная ошибка обоих методов от шага сетки (логарифмические оси)}{fig:1-2-error}
\todo{D: пояснить графики.}

\taskhead{Вывод}
\todo{D: вывод о сходимости обоих методов на треугольной области.}
